\ReviewerComment{The abstract and conclusions repeatedly refer to ``satellite-ground scheduling'' without explicitly qualifying that the results apply only to a soft-constrained, clipped-resource benchmark. Real missions impose hard energy/thermal limits. Please reword the abstract and conclusion to clearly delimit the scope, and add a short paragraph in the Discussion assessing how hard constraints might affect the observed ordering.}
\AuthorResponse{We agree. The Abstract, Introduction, Discussion, Highlights, and Conclusions now identify the study as a synthetic benchmark with clipped, soft-constrained resources. We also explain that hard constraints change the feasible action set through action masking, task rejection, or protective shutdown and may therefore change the observed ordering.}
\RevisedExcerpt{These findings apply only to this synthetic, soft-constraint benchmark. Hard energy and thermal limits can change the feasible action set and therefore the algorithm ordering.}
\ManuscriptLocation{p.~1, Abstract and Highlights; p.~3, Introduction; p.~23, Section~5.1; p.~24, Conclusions.}

\ReviewerComment{The rationale for selecting exactly these six variants is not well motivated, and the auxiliary weight $\lambda=0.3$ was calibrated only for the centred objective then reused for other auxiliary losses. This makes the component-wise comparisons (Tables 7--8) descriptive rather than fair hyperparameter-optimised contrasts. Please either run a separate grid search for each loss, or explicitly state that these comparisons are not intended to isolate the effect of centering/bootstrapping, and tone down the corresponding claims.}
\AuthorResponse{We now motivate the six objectives as matched factual-only, full-action, centered, bootstrapped, and Double-DQN controls. We retained the prespecified $\lambda=0.3$ and explicitly characterize Tables 7--8 as fixed-coefficient descriptive comparisons. The text no longer claims that these rows independently identify the effects of centering or bootstrapping.}
\RevisedExcerpt{Because the auxiliary losses differ in scale and error structure, these are fixed-coefficient implementation comparisons. They do not factorially isolate centering or bootstrapping.}
\ManuscriptLocation{p.~11, Section~3.4.1; pp.~17--18, Tables~7--8 and accompanying Results text; p.~23, Section~5.1.}

\ReviewerComment{MPC-4 is the strongest non-DQN comparator, but the literature includes more advanced optimisation methods (e.g., MILP, dynamic programming). While the authors mention this in future work, the current comparison remains limited. I request at least a literature-based discussion of how DQN might fare against such methods, and, if feasible, add one additional lightweight optimisation baseline (e.g., a greedy look-ahead heuristic) for the nominal regime.}
\AuthorResponse{We added literature-based discussion of mixed-integer programming, dynamic programming, PPO, and SAC, and we state explicitly that the study does not establish universal DQN optimality. We also implemented an eight-step greedy rollout that enumerates the first action, completes the predicted rollout with one-step greedy actions, executes the first action, and replans. It receives the same perfect forecast as MPC. The validated nominal result and timing comparison are reported in the revised Results (pp.~20--21, Table~12) and reproducibility package.}
\RevisedExcerpt{The greedy rollout reduced mean cost relative to MPC-4 (88.79 versus 90.83) but remained above all six DQN means (77.66--78.40). Every paired DQN--greedy difference was negative, ranging from $-11.12$ to $-10.39$; all 95\% bootstrap intervals excluded zero and the largest Holm-adjusted sign-test value was 0.000011.}
\ManuscriptLocation{pp.~2--3, Introduction and Related Work; pp.~12--13, Section~3.4.4; pp.~20--21, Section~4.6 and Table~12; p.~23, Section~5.2.}

\ReviewerComment{The paper provides sensitivity on task parameters but not on the discount factor $\gamma$ or training length. Since these directly affect value learning, please add a supplementary experiment or a sensitivity plot for $\gamma\in\{0.9,0.95,0.99\}$ on the nominal regime, and report learning curves for all six objectives (not just two) to confirm stable convergence.}
\AuthorResponse{We added a nominal experiment for all six objectives at $\gamma\in\{0.90,0.95,0.97,0.99\}$ with 20 seeds and the same held-out traces. MPC-4 uses the corresponding discount. We also trained each objective to 900 episodes with evaluations at 300, 600, and 900 episodes while fixing epsilon decay at 13,440 interactions. The training figure now shows all six objectives. New inference families use separate Holm corrections. These additions appear on pp.~19--20 (Figs.~6--7).}
\RevisedExcerpt{At $\gamma=0.95$, none of the six contrasts with $\gamma=0.97$ was supported after Holm correction. Reducing $\gamma$ to 0.90 increased mean cost for every objective by 2.63--4.57. All adjusted sign-test values were at or below 0.0057.}
\ManuscriptLocation{p.~14, Section~3.4.5; pp.~19--20, Sections~4.4--4.5 and Figures~6--7.}

\ReviewerComment{The ``counterfactual'' preview is clearly defined, but the misspecification test only scales costs/deltas by 0.85/1.15; it does not test structural errors (e.g., wrong thermal dynamics). Please acknowledge this limitation and suggest how more realistic misspecification could be examined in future work.}
\AuthorResponse{We agree and now distinguish amplitude perturbation from structural error. The Discussion identifies incorrect thermal recurrence, actuator saturation, state-dependent contact loss, telemetry-calibrated alternate dynamics, and hardware-in-the-loop evaluation as future tests.}
\RevisedExcerpt{Preview misspecification used uniform $\pm15\%$ scaling, which tests amplitude error but not structural error such as an incorrect thermal recurrence, unmodeled actuator saturation, or state-dependent contact loss.}
\ManuscriptLocation{p.~22, Section~5; p.~23, Sections~5.1--5.2.}
